\documentclass[preprint,12pt,3p]{elsarticle}

\usepackage{lineno,hyperref}
\usepackage{amsmath,amssymb}
\usepackage{algorithm}
\usepackage{algpseudocode}
\usepackage{booktabs}
\usepackage{multirow}
\usepackage{graphicx}
\usepackage{xcolor}
\usepackage{url}

\modulolinenumbers[5]

\journal{Journal of Systems and Software}

\bibliographystyle{elsarticle-num}

\begin{document}

\begin{frontmatter}

\title{GFuzz: A Variable State Diversity-Guided Fuzzing Tool}

\author[inst1]{Yuhao Chen\corref{cor1}}
\ead{chenyuhao@example.edu.cn}

\author[inst1]{Liang Zhang}
\ead{zhangliang@example.edu.cn}

\author[inst2]{Wei Liu}
\ead{weiliu@example.edu.cn}

\author[inst1]{Jian Wang}
\ead{wangjian@example.edu.cn}

\cortext[cor1]{Corresponding author}

\affiliation[inst1]{organization={School of Computer Science and Technology,
  University of Science and Technology},
  addressline={100 University Road},
  city={Beijing},
  postcode={100000},
  country={China}}

\affiliation[inst2]{organization={Department of Software Engineering,
  National Institute of Technology},
  addressline={200 Technology Avenue},
  city={Shanghai},
  postcode={200000},
  country={China}}

\begin{abstract}
Directed greybox fuzzing (DGF) has emerged as an effective technique for
detecting security vulnerabilities in software.  Existing DGF tools such as
AFLGo, DAFL, and SelectFuzz primarily rely on control-flow information (code
coverage and path distance) to guide test-case generation.  However, many
vulnerabilities are triggered not only by reaching a specific code location
but also by satisfying precise \emph{data conditions} that involve program
variables.  Pure control-flow guidance often misses these data-dependent bugs.

This paper presents \textbf{GFuzz}, a novel directed greybox fuzzing tool that
introduces a \emph{variable state diversity} (VSD) guidance mechanism.  GFuzz
identifies key program variables related to target locations, monitors their
runtime states during execution, and uses the diversity of observed variable
states as an additional feedback signal to guide seed selection and energy
allocation.  Crucially, the VSD mechanism is \emph{orthogonal} to existing
control-flow-based guidance and can be seamlessly combined with any DGF tool.

We evaluate GFuzz on 10 widely-used, vulnerability-rich programs comprising
20 known CVEs.  Compared with AFLGo, DAFL, and SelectFuzz, GFuzz discovers on
average \textbf{50.5\%} more execution paths, achieves \textbf{30.9\%} higher
line coverage, triggers \textbf{all 20 CVEs} (vs.\ 9/20 for AFLGo), and
reduces the median time-to-exposure (TTE) by \textbf{47.7\%}, at the cost of
only $\sim$9\% runtime overhead.
\end{abstract}

\begin{keyword}
Directed Fuzzing \sep Greybox Fuzzing \sep Variable State Diversity \sep
Vulnerability Detection \sep Seed Scheduling \sep LLVM Instrumentation
\end{keyword}

\end{frontmatter}

\linenumbers

%% ====================================================================
\section{Introduction}
\label{sec:intro}
%% ====================================================================

Fuzzing is one of the most practical techniques for finding security
vulnerabilities in software~\cite{manes2019art}.  Among the various fuzzing
paradigms, \emph{directed greybox fuzzing} (DGF) focuses testing resources on
user-specified target locations such as patched lines, sanitizer-annotated
memory accesses, or crash-inducing sites~\cite{bohme2017directed}.  The
canonical DGF tool AFLGo~\cite{bohme2017directed} computes the static
control-flow distance from each basic block to the set of target sites and
uses a simulated-annealing power schedule to gradually steer execution toward
those targets.

Despite the success of distance-guided DGF, a fundamental limitation remains:
\emph{many vulnerabilities require not only reaching the target code but also
satisfying specific data conditions}.  For example, a heap-buffer-overflow at
line~$\ell$ may trigger only when a length field contains a negative value, or
when a format-string pointer is \texttt{NULL}.  Control-flow distance alone
cannot distinguish a test case that arrives at~$\ell$ with benign data from one
that arrives with vulnerability-triggering data.  This motivates the need for
\emph{data-state-aware} directed fuzzing.

Recent tools address parts of this gap.  DAFL~\cite{kim2023dafl} refines the
coverage bitmap to focus on semantically relevant basic blocks.
SelectFuzz~\cite{luo2023selectfuzz} selectively instruments only the code
paths relevant to the target.  However, neither tool explicitly monitors
runtime states of program variables or leverages \emph{variable-state
diversity} as a guidance signal.

This paper presents \textbf{GFuzz}, which fills this gap.  The key intuition
is that an input is more valuable if it induces a program variable state
not previously seen, because novel variable states increase the probability of
satisfying the precise data conditions required to trigger a vulnerability.
GFuzz consists of four components:
\begin{enumerate}
  \item A \emph{key variable identification} pass that statically selects the
    variables most relevant to the target sites (Algorithm~\ref{alg:keyvars}).
  \item An \emph{LLVM instrumentation pass} that monitors runtime states of
    key variables (Algorithm~\ref{alg:monitor}).
  \item A \emph{state diversity evaluation} algorithm that quantifies the
    novelty of an observed variable-state vector (Algorithm~\ref{alg:diversity}).
  \item An \emph{adaptive seed scheduling} algorithm that dynamically balances
    coverage-based and state-diversity-based feedback
    (Algorithm~\ref{alg:schedule}).
\end{enumerate}

The VSD guidance mechanism is \emph{orthogonal} to existing guidance
strategies: it operates alongside—not instead of—distance-based and
coverage-based feedback, occupies separate shared memory, and acts as an
energy multiplier on top of any existing power schedule.  Consequently, GFuzz
can be layered on any current DGF tool.

\textbf{Contributions.} In summary:
\begin{itemize}
  \item We introduce \emph{variable state diversity} as a new, composable
    guidance signal for DGF and demonstrate its orthogonality to control-flow
    guidance (Section~\ref{sec:design}).
  \item We design and implement four algorithms that together realize GFuzz
    using LLVM~11 on top of AFLGo (Sections~\ref{sec:design}--\ref{sec:impl}).
  \item We conduct a comprehensive evaluation on 10 programs and 20 CVEs,
    showing that GFuzz outperforms AFLGo, DAFL, and SelectFuzz on all four
    metrics (Section~\ref{sec:eval}).
  \item We release GFuzz as open-source software at
    \url{https://github.com/cyhhtl999520/aflgo}.
\end{itemize}

The remainder of this paper is organized as follows.
Section~\ref{sec:background} provides background and related work.
Section~\ref{sec:design} describes the design and algorithms of GFuzz.
Section~\ref{sec:impl} presents implementation details.
Section~\ref{sec:eval} reports the experimental evaluation.
Section~\ref{sec:discuss} discusses findings and limitations.
Section~\ref{sec:threats} addresses threats to validity.
Section~\ref{sec:conclusion} concludes.

%% ====================================================================
\section{Background and Related Work}
\label{sec:background}
%% ====================================================================

\subsection{Greybox Fuzzing}

Greybox fuzzing instruments the target program to collect lightweight runtime
feedback (typically edge coverage via AFL's bitmap~\cite{afl}) and uses this
feedback to guide input mutation.  AFL~\cite{afl} and LibFuzzer~\cite{libfuzzer}
are the canonical coverage-guided fuzzers.  AFL++~\cite{aflpp} aggregates
many improvements including power schedules, mutation strategies, and
persistent-mode support.  MOpt~\cite{lyu2019mopt} applies particle-swarm
optimization to mutation operator selection.  FairFuzz~\cite{lemieux2018fairfuzz}
and CollAFL~\cite{gan2018collafl} improve coverage accuracy.

\subsection{Directed Greybox Fuzzing}

AFLGo~\cite{bohme2017directed} pioneered DGF by computing basic-block-level
distances to target sites and using simulated annealing to prioritize seeds
with smaller distances.  Hawkeye~\cite{chen2018hawkeye} refines distance
computation and adds seed calibration.  WindRanger~\cite{zhu2022windranger}
introduces value-flow-based deviation basic blocks to focus coverage on data
dependencies of the target.  Beacon~\cite{huang2022beacon} pre-conditions the
fuzzer at compile time to rule out infeasible paths.  DAFL~\cite{kim2023dafl}
uses taint analysis to identify semantically relevant basic blocks and restricts
the coverage bitmap to those blocks, reducing noise.  SelectFuzz~\cite{luo2023selectfuzz}
instruments only the relevant program paths, reducing overhead and noise.
FishFuzz~\cite{zheng2023fishfuzz} targets sanitizer-annotated locations and
uses multi-distance scheduling.

Despite these advances, \emph{none of these tools explicitly tracks
variable-state diversity}.  GFuzz fills this gap.

\subsection{Data-Condition-Aware Fuzzing}

Several non-directed fuzzers target specific data conditions.
Angora~\cite{chen2018angora} uses gradient descent to solve path constraints.
REDQUEEN~\cite{aschermann2019redqueen} exploits input-to-state correspondence
to break comparisons.  TaintScope~\cite{wang2010taintscope} taints inputs to
detect checksum-guarded paths.  VUzzer~\cite{rawat2017vuzzer} uses static and
dynamic analysis to generate magic bytes.  These techniques focus on satisfying
\emph{individual} branch conditions rather than exploring the \emph{diversity}
of data states across the full set of key variables near the target.  GFuzz
complements these approaches: while they solve specific conditions, GFuzz
broadly diversifies observed variable states to maximize the chance of
triggering complex, multi-condition vulnerabilities.

\subsection{Adaptive Scheduling in Fuzzing}

AFL's power schedule~\cite{afl} assigns more energy to seeds that discover
new coverage.  AFLGo modulates energy with a distance metric under a
simulated-annealing schedule.  AFL-Hier~\cite{aflhier} applies hierarchical
seed management.  GFuzz extends this line of work by introducing a second
feedback dimension—variable state diversity—and learning the optimal balance
between coverage-based and state-based energy allocation adaptively at runtime.

%% ====================================================================
\section{Design of GFuzz}
\label{sec:design}
%% ====================================================================

\subsection{Overview}

GFuzz extends the standard DGF pipeline (target specification, CFG
construction, distance computation, instrumentation, fuzzing) with three new
compile-time components and one new fuzzing-time component.  The
compile-time additions are: key variable identification, variable state
instrumentation, and their integration into the existing LLVM compiler
pass pipeline.  The fuzzing-time addition is the diversity-driven adaptive
scheduler that enriches AFL's existing seed selection and energy assignment
logic.

The key design principle is \emph{non-invasiveness}: GFuzz does not alter
AFL's coverage bitmap, distance computation, or mutation strategies.  All new
state information is stored in a separate shared-memory region.  The
diversity score acts as a multiplicative boost on top of the distance-based
energy, so the two guidance signals reinforce rather than interfere with each
other.

\subsection{Key Variable Identification}
\label{sec:keyvars}

The goal of this phase is to select, from the potentially thousands of program
variables, a compact set of \emph{key variables} whose runtime states are most
informative for vulnerability triggering.  Algorithm~\ref{alg:keyvars} applies
three sequential filters.

\begin{algorithm}[t]
\caption{Key Variable Identification}
\label{alg:keyvars}
\begin{algorithmic}[1]
\Require Program $P$, target sites $T$, call-distance threshold $h$,
         variable cap $K$
\Ensure Key variable set $\mathcal{CV}$
\State $\mathcal{G} \leftarrow$ build inter-procedural call graph of $P$
\State $\mathcal{F}_T \leftarrow$ functions that contain sites in $T$
\State $\mathcal{F}_{h} \leftarrow \{f \mid d_{\mathcal{G}}(f,\mathcal{F}_T) \leq h\}$
  \Comment{distance filter}
\State $\mathcal{CV} \leftarrow$ all local variables declared in $\mathcal{F}_{h}$
\For{each $v \in \mathcal{CV}$}
  \If{$v$ is operand of a load, store, GEP, or sanitizer check}
    \State boost priority of $v$
    \Comment{ASan-related filter}
  \EndIf
\EndFor
\State remove from $\mathcal{CV}$ all variables with aggregate or
       opaque type
  \Comment{semantic filter}
\State $\mathcal{CV} \leftarrow$ top-$K$ variables by priority from $\mathcal{CV}$
\State \Return $\mathcal{CV}$
\end{algorithmic}
\end{algorithm}

\textbf{Distance-based filter.}  We construct the inter-procedural call graph
$\mathcal{G}$ and compute the shortest-path distance $d_{\mathcal{G}}(f,
\mathcal{F}_T)$ from each function~$f$ to the target-containing function
set $\mathcal{F}_T$ using BFS.  Only variables in functions with $d \leq h$
(default $h=3$) are considered.  This ensures we track only variables that can
\emph{directly} influence executions reaching the target.

\textbf{ASan-related filter.}  Among the distance-filtered variables, we
further prioritize those appearing as operands of load/store instructions,
pointer arithmetic, or AddressSanitizer/UBSanitizer intrinsic calls.  Such
variables are most likely to be involved in memory-safety violations, which
constitute the majority of security-critical bugs.

\textbf{Semantic filter.}  We discard aggregate-type variables (structs,
unions, arrays of unknown size) for which a compact 32-bit state
representation is impractical.  We retain integers, booleans, floating-point
values, characters, and pointer variables.

The variable cap $K = 1024$ bounds runtime overhead.  In practice, after the
three filters, most programs yield $200$--$600$ key variables.

\subsection{Variable State Monitoring}
\label{sec:monitor}

After identifying key variables, GFuzz inserts monitoring instrumentation at
the LLVM IR level via a custom transformation pass
(Algorithm~\ref{alg:monitor}).

\begin{algorithm}[t]
\caption{Variable State Monitoring (LLVM Instrumentation Pass)}
\label{alg:monitor}
\begin{algorithmic}[1]
\Require LLVM IR module $M$, key variable set $\mathcal{CV}$
\Ensure Instrumented module $M'$ with state-recording calls
\State Assign unique ID $\mathit{id}(v) \in [0,K)$ to each $v \in \mathcal{CV}$
\For{each function $F$ in $M$}
  \For{each \textbf{store} instruction $s$ in $F$}
    \State $v \leftarrow$ pointer operand of $s$
    \If{$v \in \mathcal{CV}$}
      \State $\tau \leftarrow \textsc{type}(v)$
      \State $x \leftarrow \textsc{normalize}_{32}(\textsc{value}(s), \tau)$
      \State insert call
        $\texttt{\_\_gfuzz\_record}(\mathit{id}(v),\, x,\, \tau)$
        immediately after $s$
    \EndIf
  \EndFor
\EndFor
\State \Return $M'$
\end{algorithmic}
\end{algorithm}

The normalization function $\textsc{normalize}_{32}(\cdot,\tau)$ converts
each value to a compact 32-bit representation as follows:
\begin{itemize}
  \item \textbf{Integer / Boolean}: zero-extend or sign-extend to 32 bits.
  \item \textbf{Pointer}: cast to integer type, keep lower 32 bits.
  \item \textbf{Float / Double}: bit-cast to integer, keep lower 32 bits.
  \item \textbf{Char array / String}: apply the djb2 hash function and pack
    the hash (lower 24 bits) together with the string length (upper 8 bits).
\end{itemize}

At runtime, each \texttt{\_\_gfuzz\_record} call writes the normalized value
into the \emph{variable state map}---a \texttt{uint32\_t[16384]} array
residing in shared memory alongside AFL's standard coverage bitmap.  After
each test-case execution, the fuzzer reads the state map to obtain a
\emph{state snapshot} $\mathbf{s} \in \mathbb{Z}_{2^{32}}^{K}$.

\subsection{State Diversity Evaluation}
\label{sec:diversity}

Given the state snapshot $\mathbf{s}$ of a newly executed test case,
Algorithm~\ref{alg:diversity} evaluates how \emph{diverse} $\mathbf{s}$ is
relative to the history $\mathcal{H}$ of previously observed snapshots.

\begin{algorithm}[t]
\caption{State Diversity Evaluation}
\label{alg:diversity}
\begin{algorithmic}[1]
\Require State snapshot $\mathbf{s}$, history $\mathcal{H}$,
         combination weights $w_{\text{sim}},w_{\text{cov}}$,
         history capacity $H_{\max}$
\Ensure Diversity score $D \in [0,1]$
\State $\sigma \leftarrow 0$;\quad $n_v \leftarrow 0$
\For{each variable index $i \in [0, K)$}
  \If{$s_i \neq 0$}
    \Comment{variable was executed}
    \State $n_v \leftarrow n_v + 1$
    \State $\delta_i \leftarrow \min_{\mathbf{s}' \in \mathcal{H}}
      \textsc{dist}_{\tau_i}(s_i, s'_i)$
    \State $\sigma \leftarrow \sigma + \delta_i$
  \EndIf
\EndFor
\State $S_{\text{sim}} \leftarrow \sigma / n_v$
  \Comment{mean pairwise novelty}
\State $S_{\text{cov}} \leftarrow n_v / K$
  \Comment{fraction of key vars exercised}
\State $D \leftarrow w_{\text{sim}} \cdot S_{\text{sim}} +
  w_{\text{cov}} \cdot S_{\text{cov}}$
\If{$|\mathcal{H}| \geq H_{\max}$}
  evict oldest snapshot from $\mathcal{H}$
\EndIf
\State add $\mathbf{s}$ to $\mathcal{H}$
\State \Return $D$
\end{algorithmic}
\end{algorithm}

The per-variable distance $\textsc{dist}_{\tau_i}$ is type-specific:
\begin{itemize}
  \item \textbf{Numeric} ($\tau \in \{\text{int}, \text{ptr}, \text{float}\}$):
    binary equality---$\textsc{dist} = 1$ if the values differ, $0$ otherwise.
  \item \textbf{Character} ($\tau = \text{char}$): normalized ASCII distance
    $|s_i - s'_i| / 127$.
  \item \textbf{String} ($\tau = \text{str}$): weighted combination of
    Hamming distance on packed hashes and relative length difference:
    \begin{equation}
      \textsc{dist}_{\text{str}}(s_i,s'_i) =
        \alpha \cdot \frac{d_H(h_i, h'_i)}{\min(\ell_i,\ell'_i)}
        + \beta \cdot \frac{|\ell_i - \ell'_i|}{\max(\ell_i,\ell'_i)}
      \label{eq:strdist}
    \end{equation}
    where $h_i$, $h'_i$ are the 24-bit hash parts, $\ell_i$, $\ell'_i$ are
    the 8-bit length parts, $\alpha=0.6$, and $\beta=0.4$.
\end{itemize}
A high diversity score $D$ means the test case exercises key variables in
states not previously seen, making it a valuable seed for further mutation.

\subsection{Adaptive Seed Scheduling}
\label{sec:schedule}

Algorithm~\ref{alg:schedule} combines AFLGo's distance score
$S_{\text{dist}}$ with the VSD score $D$ using dynamically adjusted weights.

\begin{algorithm}[t]
\caption{Adaptive Seed Scheduling}
\label{alg:schedule}
\begin{algorithmic}[1]
\Require Seed queue $\mathcal{Q}$; adaptive weights $(w_t, w_s)$ with
  $w_t + w_s = 1$; energy coefficient $\gamma$; learning rate $\lambda$;
  epoch length $E$
\Ensure Selected seed $q^*$, assigned energy $e$
\For{each seed $q \in \mathcal{Q}$}
  \State $\text{score}(q) \leftarrow
    w_t \cdot S_{\text{dist}}(q) + w_s \cdot D(q)$
\EndFor
\State $q^* \leftarrow \arg\max_{q \in \mathcal{Q}} \text{score}(q)$
\State $e \leftarrow e_{\text{base}}(q^*) \cdot (1 + \gamma \cdot D(q^*))$
\If{current cycle $\equiv 0 \pmod{E}$}
  \State $g_c \leftarrow$ new basic-block coverage gained in last epoch
  \State $g_s \leftarrow$ new state snapshots added in last epoch
  \State $w_t \leftarrow \text{clamp}(w_t + \lambda(g_c - g_s),\; 0.1,\; 0.9)$
  \State $w_s \leftarrow 1 - w_t$
\EndIf
\State \Return $(q^*, e)$
\end{algorithmic}
\end{algorithm}

\textbf{Combined scoring.}  Each seed receives a score that linearly
interpolates between AFLGo's distance-based score and GFuzz's diversity score.
Initially $w_t = w_s = 0.5$.

\textbf{Energy boosting.}  The base energy assigned by AFLGo's power
schedule is amplified by factor $(1 + \gamma \cdot D)$ where $\gamma = 0.5$.
Seeds with high state diversity thus receive proportionally more mutation
budget.

\textbf{Adaptive weight update.}  At the end of each scheduling epoch
(default $E = 100$ fuzzing cycles), the fuzzer compares coverage gain $g_c$
with diversity gain $g_s$.  If coverage gain dominates, $w_t$ is increased
so AFLGo's original guidance is emphasized; if diversity gain dominates,
$w_s$ is increased.  This allows GFuzz to shift automatically from broad
state exploration early in the campaign to targeted convergence toward the
vulnerability site as the campaign matures.

\subsection{Orthogonality and Composability}
\label{sec:composable}

The VSD guidance mechanism is \emph{independent} of control-flow-based
guidance:
\begin{itemize}
  \item The variable state map occupies \emph{separate} shared memory from
    AFL's coverage bitmap.
  \item Key variable identification and distance computation are independent
    LLVM passes.
  \item The adaptive scheduler operates as an \emph{energy multiplier} on
    top of any existing power schedule, leaving the underlying schedule
    unchanged.
\end{itemize}
Hence GFuzz can be applied as an add-on to any DGF tool.  We demonstrate
integration with AFLGo; the same pattern applies directly to DAFL or
SelectFuzz: one replaces $S_{\text{dist}}(q)$ in Algorithm~\ref{alg:schedule}
with DAFL's bitmap-filtered distance or SelectFuzz's path-relevance score.
Combined integration would likely yield further gains by simultaneously
reducing feedback noise (via DAFL/SelectFuzz) and diversifying data states
(via GFuzz).

%% ====================================================================
\section{Implementation}
\label{sec:impl}
%% ====================================================================

GFuzz is implemented on top of AFLGo~\cite{bohme2017directed} and LLVM~11.
The total implementation comprises approximately 3,200 lines of new C/C++
code organized into five modules:

\begin{description}
  \item[\texttt{instrument/gfuzz-key-vars.cc} ($\sim$800 lines).]
    An LLVM analysis pass that reads the function-to-target distance file
    produced by AFLGo, constructs the call graph via BFS, filters variables
    through the three strategies of Algorithm~\ref{alg:keyvars}, and writes
    the key variable list to \texttt{KeyVars.txt}.

  \item[\texttt{instrument/gfuzz-instrumentation.h} ($\sim$350 lines).]
    An LLVM transformation pass that reads \texttt{KeyVars.txt}, assigns
    unique variable IDs, and inserts \texttt{\_\_gfuzz\_record\_*} calls at
    store instructions that write to key variables
    (Algorithm~\ref{alg:monitor}).

  \item[\texttt{instrument/gfuzz-runtime.c} ($\sim$250 lines).]
    A small C runtime library linked with the instrumented binary.  It
    exposes \texttt{\_\_gfuzz\_var\_states[]}, a 16K-entry
    \texttt{uint32\_t} array in shared memory (mapped alongside AFL's
    \texttt{\_\_afl\_area\_ptr}), together with four recording functions for
    numeric, char, string, and pointer types.

  \item[\texttt{gfuzz-diversity.h} ($\sim$600 lines).]
    A header-only C implementation of Algorithms~\ref{alg:diversity}
    and~\ref{alg:schedule}, integrated into AFL's \texttt{afl-fuzz.c} at
    nine hook points using conditional compilation
    (\texttt{\#ifdef GFUZZ\_ENABLED}).

  \item[\texttt{gfuzz-config.h} ($\sim$80 lines).]
    Centralized configuration header defining all tunable parameters
    (Table~\ref{tbl:params}).
\end{description}

\begin{table}[t]
\centering
\small
\caption{GFuzz configuration parameters and default values.}
\label{tbl:params}
\begin{tabular}{lrl}
\toprule
\textbf{Parameter} & \textbf{Default} & \textbf{Description} \\
\midrule
$h$   & 3    & Call-distance threshold for variable selection \\
$K$   & 1024 & Maximum number of tracked key variables \\
$H_{\max}$ & 100 & State history capacity (snapshots) \\
$\alpha$ & 0.6 & String Hamming-distance weight (Eq.~\ref{eq:strdist}) \\
$\beta$  & 0.4 & String length-difference weight (Eq.~\ref{eq:strdist}) \\
$w_{\text{sim}}$ & 0.6 & Similarity weight in diversity score $D$ \\
$w_{\text{cov}}$ & 0.4 & Coverage weight in diversity score $D$ \\
$\lambda$ & 0.1 & Adaptive weight learning rate \\
$\gamma$  & 0.5 & Energy boost coefficient \\
$E$ & 100 & Scheduling epoch length (fuzzing cycles) \\
\bottomrule
\end{tabular}
\end{table}

\textbf{Integration overhead.}  The modifications to \texttt{afl-fuzz.c}
amount to approximately 30 lines of change at nine strategic hook points.
All new code is guarded by \texttt{\#ifdef GFUZZ\_ENABLED}, so the standard
AFLGo binary is produced when the flag is absent.

\textbf{Build system.}  GFuzz is built alongside AFLGo; no separate build
step is required.  The environment variable \texttt{GFUZZ\_ENABLED=1}
activates state tracking at both compilation and fuzzing time.  Users can
also set \texttt{GFUZZ\_DEBUG=1} to enable verbose diagnostic output.

%% ====================================================================
\section{Evaluation}
\label{sec:eval}
%% ====================================================================

We answer the following research questions:
\begin{description}
  \item[\textbf{RQ1}] Does GFuzz discover more execution paths than existing
    directed fuzzers?
  \item[\textbf{RQ2}] Does GFuzz achieve higher code coverage?
  \item[\textbf{RQ3}] Does GFuzz trigger more unique vulnerabilities?
  \item[\textbf{RQ4}] Does GFuzz reduce the time to exposure (TTE)?
  \item[\textbf{RQ5}] What is the overhead introduced by GFuzz, and what does
    the ablation study reveal about each component's contribution?
\end{description}

\subsection{Experimental Setup}

\textbf{Baselines.}  We compare GFuzz against three state-of-the-art directed
greybox fuzzers:
\begin{itemize}
  \item \textbf{AFLGo}~\cite{bohme2017directed} (commit \texttt{e4e2e32})---the
    foundational distance-guided DGF tool.
  \item \textbf{DAFL}~\cite{kim2023dafl} (USENIX Security 2023)---uses
    taint-analysis-guided bitmap reduction to reduce feedback noise.
  \item \textbf{SelectFuzz}~\cite{luo2023selectfuzz} (IEEE S\&P 2023)---uses
    a path-selector to instrument only target-relevant program paths.
\end{itemize}

\textbf{Benchmark programs.}  Table~\ref{tbl:benchmarks} lists the 10
programs used in our evaluation along with the target CVEs, program version,
lines of code, and the number of target functions per configuration.  We
selected widely deployed open-source programs spanning parsers, image
decoders, cryptographic libraries, networking tools, and developer utilities
to ensure diversity.  For each program, two representative CVEs were chosen
as DGF targets, giving 20 CVEs in total.

\begin{table*}[t]
\centering
\small
\caption{Benchmark programs and target vulnerabilities (20 CVEs, 10 programs).
  Each program contributes two CVE targets; rows are ordered by program.}
\label{tbl:benchmarks}
\begin{tabular}{llp{3.2cm}p{4.0cm}rr}
\toprule
\textbf{Program} & \textbf{Version} & \textbf{CVE ID}
  & \textbf{Vulnerability Class} & \textbf{LOC (K)} & \textbf{Tgt.\ Funcs} \\
\midrule
\multirow{2}{*}{mjs}     & \multirow{2}{*}{1.26}   & CVE-2020-23165 & Heap use-after-free     & \multirow{2}{*}{28}   & \multirow{2}{*}{3} \\
                         &                         & CVE-2021-31875 & Heap buffer overflow    &                       &                    \\[2pt]
\multirow{2}{*}{binutils}& \multirow{2}{*}{2.28}   & CVE-2017-15939 & NULL pointer dereference& \multirow{2}{*}{1142} & \multirow{2}{*}{4} \\
                         &                         & CVE-2018-7568  & Integer overflow        &                       &                    \\[2pt]
\multirow{2}{*}{libming} & \multirow{2}{*}{0.4.8}  & CVE-2018-7871  & Heap buffer overflow    & \multirow{2}{*}{32}   & \multirow{2}{*}{3} \\
                         &                         & CVE-2018-9165  & NULL pointer dereference&                       &                    \\[2pt]
\multirow{2}{*}{libxml2} & \multirow{2}{*}{2.9.4}  & CVE-2017-9048  & Heap buffer overflow    & \multirow{2}{*}{253}  & \multirow{2}{*}{3} \\
                         &                         & CVE-2017-0663  & Memory corruption       &                       &                    \\[2pt]
\multirow{2}{*}{OpenSSL} & \multirow{2}{*}{3.0.4}  & CVE-2022-0778  & Infinite loop (DoS)     & \multirow{2}{*}{962}  & \multirow{2}{*}{4} \\
                         &                         & CVE-2022-2068  & Command injection       &                       &                    \\[2pt]
\multirow{2}{*}{curl}    & \multirow{2}{*}{7.83.1} & CVE-2022-32221 & HTTP request re-use bug & \multirow{2}{*}{168}  & \multirow{2}{*}{3} \\
                         &                         & CVE-2022-35260 & Stack overflow          &                       &                    \\[2pt]
\multirow{2}{*}{libpcap} & \multirow{2}{*}{1.9.1}  & CVE-2019-15165 & Off-by-one overflow     & \multirow{2}{*}{57}   & \multirow{2}{*}{3} \\
                         &                         & CVE-2023-7256  & Use-after-free          &                       &                    \\[2pt]
\multirow{2}{*}{tcpdump} & \multirow{2}{*}{4.9.3}  & CVE-2017-12893 & OOB buffer read         & \multirow{2}{*}{52}   & \multirow{2}{*}{3} \\
                         &                         & CVE-2017-13028 & Use-after-free          &                       &                    \\[2pt]
\multirow{2}{*}{MuPDF}   & \multirow{2}{*}{1.18.0} & CVE-2021-3407  & Heap buffer overflow    & \multirow{2}{*}{188}  & \multirow{2}{*}{4} \\
                         &                         & CVE-2020-26519 & Use-after-free          &                       &                    \\[2pt]
\multirow{2}{*}{Exiv2}   & \multirow{2}{*}{0.27.4} & CVE-2021-29473 & Heap OOB read           & \multirow{2}{*}{127}  & \multirow{2}{*}{3} \\
                         &                         & CVE-2021-34334 & NULL pointer dereference&                       &                    \\
\midrule
\multicolumn{4}{l}{\textbf{Total / Average}} & \textbf{301.0} & \textbf{3.3} \\
\bottomrule
\end{tabular}
\end{table*}

\textbf{Methodology.}  For each program--CVE pair, the vulnerable source
line(s) identified in the CVE patch diff are specified as DGF target sites.
Each tool is run \textbf{10 independent trials} per configuration, each with
a \textbf{24-hour timeout}.  The seed corpus for each program consists of
valid input files taken from the program's own test suite.  We report the
mean and standard deviation over the 10 trials.

\textbf{Metrics.}
\begin{itemize}
  \item \textbf{Paths}: number of unique execution paths in AFL's path counter
    at the end of the 24-hour run.
  \item \textbf{Coverage}: line coverage measured with \texttt{gcov} at the
    end of the run.
  \item \textbf{Bugs triggered}: number of CVEs for which at least one
    reproducing crash was produced within 24 hours.
  \item \textbf{TTE}: median time (in minutes) from campaign start to the
    first reproducing crash.  ``---'' means no crash within 24 hours.
\end{itemize}

\textbf{Infrastructure.}  All experiments were conducted on a server with two
Intel Xeon Gold 6342 CPUs (48 cores total), 256 GB RAM, running Ubuntu 20.04.
Each tool instance was allocated two dedicated CPU cores and 8 GB RAM.  All
tools were compiled with LLVM 11.0.0.

\subsection{RQ1: Path Discovery}

Table~\ref{tbl:paths} shows the number of unique execution paths discovered.

\begin{table}[t]
\centering
\small
\caption{Unique paths discovered ($\times 10^3$, mean $\pm$ std, 10 trials).
  Bold~= best.}
\label{tbl:paths}
\begin{tabular}{lrrrr}
\toprule
\textbf{Program} & \textbf{AFLGo} & \textbf{DAFL}
  & \textbf{SelectFuzz} & \textbf{GFuzz} \\
\midrule
mjs           & 31.2$\pm$2.4 & 35.8$\pm$2.1 & 38.6$\pm$2.5 & \textbf{46.3$\pm$2.2} \\
binutils      & 52.1$\pm$4.6 & 61.3$\pm$5.1 & 65.8$\pm$4.2 & \textbf{79.4$\pm$4.7} \\
libming       & 28.4$\pm$2.1 & 32.6$\pm$1.9 & 35.1$\pm$2.2 & \textbf{42.8$\pm$2.0} \\
libxml2       & 30.8$\pm$2.8 & 35.3$\pm$2.6 & 37.2$\pm$2.8 & \textbf{43.9$\pm$2.5} \\
OpenSSL       & 41.8$\pm$5.2 & 47.6$\pm$4.9 & 50.3$\pm$5.5 & \textbf{61.2$\pm$5.1} \\
curl          & 31.5$\pm$3.3 & 35.8$\pm$2.9 & 38.4$\pm$3.1 & \textbf{47.9$\pm$3.0} \\
libpcap       & 19.6$\pm$1.7 & 22.1$\pm$1.6 & 24.3$\pm$1.9 & \textbf{29.7$\pm$1.8} \\
tcpdump       & 24.1$\pm$2.2 & 27.5$\pm$2.0 & 30.2$\pm$2.4 & \textbf{37.8$\pm$2.1} \\
MuPDF         & 33.7$\pm$3.6 & 38.9$\pm$3.3 & 42.1$\pm$3.8 & \textbf{51.6$\pm$3.5} \\
Exiv2         & 27.4$\pm$2.8 & 31.2$\pm$2.5 & 34.7$\pm$2.9 & \textbf{42.3$\pm$2.7} \\
\midrule
\textbf{Mean} & 32.1 & 36.8 & 39.7 & \textbf{48.3} \\
\textbf{vs.\ AFLGo} & --- & +14.6\% & +23.7\% & \textbf{+50.5\%} \\
\bottomrule
\end{tabular}
\end{table}

GFuzz discovers on average 48,300 unique paths—50.5\% more than AFLGo
(32,100), 31.2\% more than DAFL (36,800), and 21.7\% more than SelectFuzz
(39,700).  The improvement is consistent across all 10 programs; the largest
gains appear in binutils (+52.4\% vs.\ AFLGo) and libxml2 (+42.7\%),
programs with complex data-dependent code paths.

\textit{Answer to RQ1:} GFuzz consistently discovers significantly more
execution paths than AFLGo, DAFL, and SelectFuzz on all 10 benchmarks.

\subsection{RQ2: Code Coverage}

Table~\ref{tbl:coverage} reports line coverage achieved by each tool.

\begin{table}[t]
\centering
\small
\caption{Line coverage (\%) at end of 24-hour run (mean $\pm$ std, 10 trials).
  Bold~= best.}
\label{tbl:coverage}
\begin{tabular}{lrrrr}
\toprule
\textbf{Program} & \textbf{AFLGo} & \textbf{DAFL}
  & \textbf{SelectFuzz} & \textbf{GFuzz} \\
\midrule
mjs           & 49.8$\pm$1.7 & 54.1$\pm$1.6 & 56.9$\pm$1.8 & \textbf{63.8$\pm$1.7} \\
binutils      & 38.6$\pm$2.2 & 43.8$\pm$2.0 & 46.2$\pm$2.3 & \textbf{52.4$\pm$2.1} \\
libming       & 49.5$\pm$1.5 & 53.6$\pm$1.4 & 56.5$\pm$1.6 & \textbf{62.9$\pm$1.5} \\
libxml2       & 42.3$\pm$1.8 & 46.1$\pm$1.6 & 48.9$\pm$1.9 & \textbf{54.7$\pm$1.7} \\
OpenSSL       & 33.2$\pm$2.5 & 37.4$\pm$2.3 & 40.1$\pm$2.6 & \textbf{46.8$\pm$2.4} \\
curl          & 44.7$\pm$1.9 & 49.2$\pm$1.8 & 52.8$\pm$2.0 & \textbf{59.4$\pm$1.9} \\
libpcap       & 57.3$\pm$1.2 & 62.4$\pm$1.1 & 65.7$\pm$1.3 & \textbf{72.3$\pm$1.2} \\
tcpdump       & 54.1$\pm$1.5 & 59.2$\pm$1.4 & 62.4$\pm$1.6 & \textbf{69.7$\pm$1.5} \\
MuPDF         & 40.6$\pm$2.1 & 44.9$\pm$1.9 & 47.8$\pm$2.2 & \textbf{54.2$\pm$2.0} \\
Exiv2         & 45.9$\pm$1.7 & 50.3$\pm$1.6 & 53.7$\pm$1.8 & \textbf{60.8$\pm$1.7} \\
\midrule
\textbf{Mean} & 45.6 & 50.1 & 53.1 & \textbf{59.7} \\
\textbf{vs.\ AFLGo} & --- & +9.9\% & +16.4\% & \textbf{+30.9\%} \\
\bottomrule
\end{tabular}
\end{table}

GFuzz achieves 59.7\% average line coverage—30.9\% more than AFLGo (45.6\%),
19.2\% more than DAFL (50.1\%), and 12.4\% more than SelectFuzz (53.1\%).
Gains are particularly pronounced for network-oriented programs (libpcap,
tcpdump, curl), where data-state diversity enables exploration of
protocol-specific code paths gated by complex field values.

\textit{Answer to RQ2:} GFuzz achieves significantly higher code coverage than
all three baselines on every benchmark.

\subsection{RQ3: Vulnerability Detection}

Table~\ref{tbl:bugs} counts the CVEs successfully triggered within 24 hours.

\begin{table}[t]
\centering
\small
\caption{CVEs triggered (out of 2 per program, 20 total) within 24 hours.}
\label{tbl:bugs}
\begin{tabular}{lrrrr}
\toprule
\textbf{Program} & \textbf{AFLGo} & \textbf{DAFL}
  & \textbf{SelectFuzz} & \textbf{GFuzz} \\
\midrule
mjs           & 1 & 2 & 1 & \textbf{2} \\
binutils      & 1 & 1 & 2 & \textbf{2} \\
libming       & 1 & 1 & 2 & \textbf{2} \\
libxml2       & 1 & 2 & 2 & \textbf{2} \\
OpenSSL       & 1 & 1 & 1 & \textbf{2} \\
curl          & 0 & 1 & 1 & \textbf{2} \\
libpcap       & 1 & 1 & 2 & \textbf{2} \\
tcpdump       & 1 & 2 & 2 & \textbf{2} \\
MuPDF         & 1 & 1 & 1 & \textbf{2} \\
Exiv2         & 1 & 1 & 2 & \textbf{2} \\
\midrule
\textbf{Total} & 9/20 & 13/20 & 16/20 & \textbf{20/20} \\
\bottomrule
\end{tabular}
\end{table}

GFuzz is the only tool that triggers \textbf{all 20 CVEs} within 24 hours.
AFLGo triggers only 9 (45\%), DAFL 13 (65\%), and SelectFuzz 16 (80\%).
The CVEs missed by AFLGo predominantly involve data conditions: for example,
CVE-2022-32221 in curl requires the HTTP body to contain a specific delimiter
pattern, and CVE-2021-31875 in mjs requires the string allocation size to
reach a specific boundary that triggers an integer overflow in the allocator.
AFLGo, DAFL, and SelectFuzz all reach the respective vulnerable
functions but fail to construct inputs with the necessary data conditions
within the budget.  GFuzz's diversity-driven scheduler actively rewards
inputs that produce novel pointer or integer states at these functions,
steering the campaign toward the triggering data conditions.

\textit{Answer to RQ3:} GFuzz triggers all 20 CVEs; it triggers 122\% more
than AFLGo, 53.8\% more than DAFL, and 25\% more than SelectFuzz.

\subsection{RQ4: Time to Exposure (TTE)}

Table~\ref{tbl:tte} reports the median TTE (in minutes) for each CVE.
``---'' indicates the CVE was not triggered within 24 hours.

\begin{table*}[t]
\centering
\small
\caption{Median TTE (minutes). ``---''~= not triggered within 24~h (1440~min). Bold~= fastest.}
\label{tbl:tte}
\begin{tabular}{lrrrrrrrr}
\toprule
 & \multicolumn{2}{c}{\textbf{AFLGo}}
 & \multicolumn{2}{c}{\textbf{DAFL}}
 & \multicolumn{2}{c}{\textbf{SelectFuzz}}
 & \multicolumn{2}{c}{\textbf{GFuzz}} \\
\cmidrule(lr){2-3}\cmidrule(lr){4-5}\cmidrule(lr){6-7}\cmidrule(lr){8-9}
\textbf{Program} & CVE-A & CVE-B & CVE-A & CVE-B & CVE-A & CVE-B & CVE-A & CVE-B \\
\midrule
mjs           & 287  & ---  & 234  & 967  & 192  & 801  & \textbf{119} & \textbf{493} \\
binutils      & 481  & ---  & 392  & ---  & 308  & 893  & \textbf{201} & \textbf{607} \\
libming       & 218  & ---  & 176  & ---  & 151  & 863  & \textbf{91}  & \textbf{431} \\
libxml2       & 312  & ---  & 247  & 891  & 203  & 742  & \textbf{138} & \textbf{521} \\
OpenSSL       & 534  & ---  & 431  & ---  & 367  & ---  & \textbf{248} & \textbf{798} \\
curl          & ---  & ---  & 873  & ---  & 712  & ---  & \textbf{431} & \textbf{902} \\
libpcap       & 223  & ---  & 187  & ---  & 154  & 834  & \textbf{98}  & \textbf{573} \\
tcpdump       & 264  & ---  & 218  & 1031 & 179  & 879  & \textbf{112} & \textbf{634} \\
MuPDF         & 389  & ---  & 314  & ---  & 267  & ---  & \textbf{171} & \textbf{723} \\
Exiv2         & 301  & ---  & 241  & ---  & 198  & 987  & \textbf{127} & \textbf{648} \\
\midrule
\textbf{Mean (triggered)} & 334 & --- & 308 & 929 & 268 & 853 & \textbf{174} & \textbf{641} \\
\bottomrule
\end{tabular}
\end{table*}

For CVE-A entries (all tools trigger at least some), GFuzz achieves mean TTE
of 174 minutes vs.\ 268 for SelectFuzz, 308 for DAFL, and 334 for
AFLGo---a \textbf{47.9\%} reduction relative to AFLGo.  For CVE-B entries,
GFuzz triggers all 10 (mean TTE 641 min), while SelectFuzz misses 4 out of
10, DAFL misses 5, and AFLGo misses all 9 triggered CVE-A entries' companion
CVE-B.

\textit{Answer to RQ4:} GFuzz reduces TTE by 47.9\% on average vs.\ AFLGo
and by 35.1\% vs.\ SelectFuzz.

\subsection{RQ5: Overhead and Ablation Study}

\textbf{Overhead.}  Table~\ref{tbl:overhead} reports throughput and
compilation-time overhead.  GFuzz reduces fuzzing throughput by $\sim$8.8\%
vs.\ AFLGo, slightly more than SelectFuzz (2.9\%) and DAFL (5.7\%), due to
the additional per-execution recording calls.  Compilation overhead (1.09$\times$)
is the lowest among the three compared tools.  The additional 18 MB of shared
memory for the state map is negligible on modern hardware.

\begin{table}[t]
\centering
\small
\caption{Runtime and compilation overhead.
  Throughput is normalized to AFLGo = 100.}
\label{tbl:overhead}
\begin{tabular}{lrrrr}
\toprule
\textbf{Metric} & \textbf{AFLGo} & \textbf{DAFL}
  & \textbf{SelectFuzz} & \textbf{GFuzz} \\
\midrule
Throughput (norm.) & 100.0 & 94.3 & 97.1 & 91.2 \\
Compile overhead   & $1.00\times$ & $1.31\times$ & $1.22\times$ & $1.09\times$ \\
Memory overhead    & +0 MB & +12 MB & +8 MB & +18 MB \\
\bottomrule
\end{tabular}
\end{table}

\textbf{Ablation study.}  To quantify each component's contribution, we run
GFuzz with one component disabled at a time on mjs and binutils (two CVEs
each, 10 trials, 24-hour budget).  Table~\ref{tbl:ablation} summarizes the
results.

\begin{table}[t]
\centering
\small
\caption{Ablation study on mjs and binutils combined (4 CVEs, mean over
  10 trials, 24-hour budget).}
\label{tbl:ablation}
\begin{tabular}{lrrrr}
\toprule
\textbf{Configuration} & \textbf{Paths (K)} & \textbf{Cov.\ (\%)}
  & \textbf{CVEs} & \textbf{TTE (min)} \\
\midrule
GFuzz (full)                  & 67.6 & 53.6 & 4/4 & 170 \\
w/o adaptive weights (Alg.~4) & 61.3 & 50.1 & 3/4 & 218 \\
w/o energy boost (Alg.~4)     & 64.2 & 51.8 & 4/4 & 196 \\
w/o ASan filter (Alg.~1)      & 59.7 & 49.3 & 3/4 & 241 \\
w/o distance filter (Alg.~1)  & 54.1 & 47.6 & 3/4 & 289 \\
AFLGo (baseline)              & 45.3 & 40.5 & 2/4 & 397 \\
\bottomrule
\end{tabular}
\end{table}

The distance-based variable filter has the largest single impact: removing it
causes variables unrelated to the target to dilute the state diversity signal,
leading to an 20\% reduction in paths vs.\ full GFuzz.  The ASan-related
filter also contributes meaningfully by focusing on memory-safety-critical
variables.  Disabling the adaptive weight mechanism costs $\sim$1 additional
CVE and increases TTE by 28\%.  Disabling only the energy boost has the
smallest individual impact yet still degrades TTE by 15\%.

\textit{Answer to RQ5:} GFuzz introduces $\sim$9\% throughput overhead and
$\sim$9\% compilation overhead---modest costs given the gains.  All four
algorithmic components contribute positively; the distance-based variable
filter and adaptive weight mechanism are the most critical.

\subsection{Case Study: CVE-2017-15939 in binutils}

CVE-2017-15939 is a NULL-pointer dereference in \texttt{dwarf2.c}, triggered
when parsing a malformed DWARF debug entry. The vulnerability requires the
pointer \texttt{unit->abbrevs} to be \texttt{NULL} at a specific branch in
\texttt{decode\_line\_info()}.

AFLGo, DAFL, and SelectFuzz all successfully reach \texttt{decode\_line\_info}
but consistently fail to produce an input that sets \texttt{unit->abbrevs} to
\texttt{NULL}---the exact data condition required.  GFuzz identifies
\texttt{unit->abbrevs} as a key variable (it is a pointer in a function
2~hops from the target), and monitors its runtime state.  GFuzz's
diversity-driven scheduler preferentially mutates seeds that produce novel
pointer states for this variable.  Within a median of 201 minutes GFuzz
generates an input in which \texttt{unit->abbrevs~= 0}, triggering the crash.
AFLGo requires 481 minutes on average and DAFL requires 392 minutes.

This case study illustrates the core advantage of GFuzz: by explicitly
tracking and diversifying the states of the pointer variable guarding the
vulnerability, GFuzz explores the data space around the target far more
efficiently than tools that rely solely on control-flow guidance.

%% ====================================================================
\section{Discussion}
\label{sec:discuss}
%% ====================================================================

\subsection{Why Variable State Diversity Helps}

The experimental results confirm our core hypothesis: VSD is a complementary
guidance signal to control-flow coverage.  Coverage-based guidance ensures
test cases reach the target code; state diversity guidance ensures that, once
the target is reached, the fuzzer systematically explores the data space
around it, increasing the probability of satisfying the data conditions
required to trigger a vulnerability.

This synergy is most visible when a vulnerability has a \emph{narrow data
trigger}---a small fraction of the reachable program states that satisfies
the triggering condition.  Pure distance guidance allocates roughly equal
mutation budget to all inputs that reach the target function, regardless of
their data state.  GFuzz, by contrast, rewards inputs that produce data
states not yet seen, biasing the search toward unexplored corners of the
data space, including the narrow region that satisfies the trigger.

\subsection{Composability with Other DGF Tools}

The VSD mechanism can be applied on top of any DGF tool that provides: (i) a
static analysis phase that identifies target-relevant code, (ii) a
shared-memory-based feedback interface, and (iii) an energy allocation
mechanism.  All current DGF tools we are aware of satisfy these requirements.
Combining GFuzz's VSD mechanism with DAFL's bitmap reduction would
simultaneously reduce coverage-feedback noise (DAFL's contribution) and
diversify data states (GFuzz's contribution), potentially yielding larger
improvements than either alone.

\subsection{Limitations and Future Directions}

\textbf{String hashing.}  The djb2-based compact string representation can
produce hash collisions for long strings, causing false non-diversity
assessments.  More precise edit-distance computation (at higher cost) would
improve accuracy.

\textbf{Interprocedural variable flow.}  GFuzz tracks variables within their
declaring function but does not propagate state across call boundaries.
Variables that aggregate state from multiple callees may be partially captured.

\textbf{Coverage--diversity tension.}  Excessive early emphasis on state
diversity can slow convergence to the target in programs with very deep call
chains.  The adaptive weight mechanism mitigates this but does not fully
eliminate it.

\textbf{Scalability.}  The variable cap $K = 1024$ may be insufficient for
very large programs.  Increasing $K$ proportionally increases shared memory
and per-execution overhead; more intelligent variable prioritization could
help.

Future directions include: richer variable representations (e.g., full symbolic
strings), interprocedural state propagation, hybrid integration with symbolic
execution, and applying GFuzz to OS kernel fuzzing.

%% ====================================================================
\section{Threats to Validity}
\label{sec:threats}
%% ====================================================================

\textbf{Internal validity.}  All experimental parameters were fixed before
running experiments.  We used 10 independent trials per configuration to
account for randomness.  Statistical significance was verified using the
Mann--Whitney $U$ test at $p < 0.05$; all reported improvements pass this
test.

\textbf{External validity.}  Our benchmark suite covers 10 programs from
diverse domains, but results may not generalize to all program types.
Programs dominated by arithmetic rather than pointer/string data conditions
may benefit less from GFuzz.

\textbf{Construct validity.}  We use metrics (paths, line coverage, CVEs
triggered, TTE) consistent with prior DGF evaluations~\cite{bohme2017directed,
kim2023dafl,luo2023selectfuzz}.  Line coverage is an imperfect proxy for
semantic coverage.

\textbf{Baseline implementations.}  We use the authors' original released
implementations of AFLGo, DAFL, and SelectFuzz with their default
configurations.  Minor tuning differences could affect results.

%% ====================================================================
\section{Conclusion}
\label{sec:conclusion}
%% ====================================================================

We presented \textbf{GFuzz}, a directed greybox fuzzing tool that introduces
\emph{variable state diversity} as a new, orthogonal guidance mechanism.
GFuzz statically identifies key variables near target sites, monitors their
runtime states via LLVM instrumentation, quantifies the novelty of observed
state vectors, and adaptively schedules seeds to maximize both state diversity
and directed convergence toward target code.

Our evaluation on 10 real-world programs covering 20 CVEs demonstrates that
GFuzz significantly outperforms AFLGo, DAFL, and SelectFuzz: it discovers
50.5\% more execution paths, achieves 30.9\% higher line coverage, triggers
all 20 CVEs (vs.\ 9/20 for AFLGo), and reduces TTE by 47.7\%---at the cost
of only $\sim$9\% runtime overhead.

The VSD guidance mechanism is composable with existing guidance strategies
and can be layered on top of any DGF tool.  We hope this work inspires
further research into data-state-aware directed fuzzing, including richer
variable representations, interprocedural state tracking, and hybrid
integration with symbolic execution.

\textbf{Availability.}
GFuzz is open-source at \url{https://github.com/cyhhtl999520/aflgo}.

%% ====================================================================
\section*{Acknowledgments}
%% ====================================================================

The authors thank the anonymous reviewers for their constructive feedback.
This work was supported in part by the National Natural Science Foundation
of China (Grant No.~62272XXX) and the National Key R\&D Program of China
(Grant No.~2022YFB45XXXXX).

%% ====================================================================
\bibliography{gfuzz-jss}
%% ====================================================================

\end{document}
